\section{A nonzero downward lift with a varying Schwartz function}

% BEGIN Y87R1
\paragraph{Y87R1. A Schwartz function detecting the original principal level.}
Let $A$ be the positive definite source of a downward lift and let
$W$ be its target. In the notation of the source's $\mathrm{DownNe}$,
put $r=\dim_E W$, $n=\dim_E A=r+t$, and $S=\mathrm{under}(Z)$.
Use the given frame $\beta$, Haar measure, additive character and
splitting characters. Suppose the hypotheses $\mathrm{Defin}$,
$\mathrm{Places}$, $\mathrm{Form}$ and $\mathrm{Occ}$ hold. In
particular $0<r$ and $t\ge r+5$. Write $K_f$ for the finite level
which is the identity modulo each selected prime of $Z$ in this
frame and is the lattice stabilizer away from $S$.

Adjoin $t$ hyperbolic planes to $W$, writing $H$ for their sum and
$B=W\perp H$. Use the trivial character of the $H$-space in the
splitting on $U(A)$. There is a nonzero Schwartz function $f$ for
$(A,H)$ which is fixed by $U(A)_\infty K_f$, spherical away from
$S$, and whose Weil matrix coefficient
$\langle\omega_H(g) f,f\rangle$ vanishes for every $g$ with
integral components outside $S$ and finite component outside $K_f$.

Indeed, the support idempotent $e$ for $Z$ satisfies $e\bar e=0$.
Thus each place in $S$ is split and exactly one of the two places
above it belongs to $Z$. By W87N1 the action of $U(A)$ on the auxiliary
space is the specified Schr\"odinger change of variables. At such a place use the
selected factor of the symplectic space as a polarization. In frames the variable is
$X\in\operatorname{Mat}_{2t,n}(F_v)$. This polarization comes from
the two factors of $E_v=F_v\times F_v$. It retains all $2t$
coordinates of $H$ on the selected side. Since $2t\ge n$, the
following compact open set is nonempty:
\[
 \mathcal U_v=\left\{X\in\operatorname{Mat}_{2t,n}(\mathcal O_v):
 X\equiv\binom{I_n}{0}\pmod{\mathfrak p_v}\right\}.
\]
Use its characteristic function. The source action on this model
is right change of variables by $g_v^{\pm1}$ with the determinant
modulus of the Schr\"odinger action. Its splitting character on the
source here is trivial. If $X$ and $Xg_v$ both lie in $\mathcal U_v$,
their top square blocks lie in $I_n+\mathfrak p_v M_n(\mathcal O_v)$
and $g_v$ is the quotient of those two blocks. Hence it is in the
same principal subgroup. On that subgroup the set is preserved and
the modulus equals one. For $g_v$ outside it the supports of the two
functions are disjoint, so their inner product is zero. The argument
also applies to the inverse change of variables. In the oriented
lattice model, the characteristic function is transported from the
split Schr\"odinger model by the theta transform. The transform
preserves orthogonality up to a positive measure factor by Fourier
isometry. Equivalently, the characteristic function of its support
on a common finite lattice quotient is a sum of additive characters.
The corresponding Heisenberg projection fixes the first function
and kills its translate when the supports are disjoint.

At the remaining finite places use the self-dual lattice vacuum.
It is stabilized on the source side by the integral unitary group:
in the hyperbolic polarization of $H$ the action on $t$ copies of
$A$ is change of variables, the integral lattice is preserved and
the source splitting character is one. At infinity use the Gaussian
for the positive metric on $A$ and the balanced polarization of
$H$. The source acts unitarily there. Its two vacuum determinant
weights from the $(t,t)$ signature cancel, with no remaining
splitting character. The product of these choices is the desired
function, with nonzero norm for every covolume-one normalization.
% END Y87R1

% BEGIN Y87R2
\paragraph{Y87R2. The factor contributed by the excluded places and infinity.}
Let $u,\Phi$ be the scalar cusp form and the polynomial Fock test in
$\mathrm{Form}$ and $\mathrm{Occ}$. On $(A,B)$ take
$\Psi=\Phi\boxtimes f$ using Y87R1. Put $b=\dim_E B=r+2t$ and
$s_*=(b-n)/2=t/2$. On the doubled group $U(A\oplus-A)$ start with
the Siegel--Weil section for $\delta(\Psi,\Psi)$ at $s_*$, where
$\delta$ is the tensor with its contragredient in the second copy.
Extend it to the standard flat section by fixing its restriction to
the given maximal compact. In a Schr\"odinger model the original
section at $g$ is evaluation of the transformed test at zero.
By W87N2, in the oriented model evaluation is the integral along the diagonal
Lagrangian modulo rational points. The partial Fourier transform in
that Lagrangian gives the same evaluation: Fourier inversion uses
the character of the given Weil representation and the
covolume-one measure. Thus on the embedded $(g,1)$ the flat section is
\[
  \langle\omega_B(g)\Psi,\Psi\rangle
                  H(g)^{s-s_*},
\]
where $H$ is the Siegel height from Iwasawa decomposition and the
inner product is linear in the translated function.

Restrict the unfolded integrand to $g$ integral outside $S$. The
orthogonal sum of targets factors its Weil coefficient into the
$W$-coefficient of $\Phi$ and the $H$-coefficient of $f$: on
functions this is an external tensor, and the two covolume-one
integration measures give the product measure. Y87R1 restricts the
remaining integral to $U(A)_\infty K_f$. On this subgroup the
embedded element $(g,1)$ lies in the chosen doubled maximal compact.
At a finite place this follows by preservation of the frame lattice
in each summand. At infinity the original definite frame identifies
the metric with a positive unitary metric, so every $g_\infty$
preserves the direct sum metric. Hence $H(g)=1$ throughout this
integral. All three vectors $u,\Phi,f$ are fixed by $K_f$. Pushing
the restricted Haar measure forward to $U(A)_\infty$ therefore
gives a Haar measure $dm_\infty$ there, with the positive volume of
$K_f$ included. Since $f$ is fixed at infinity, the integral equals
\[
  C=\|f\|^2\int_{U(A)_\infty}
   \langle\omega_W(k)\Phi,\Phi\rangle\,
   \langle R(k)u,u\rangle\,dm_\infty(k).
\]
This is nonzero by the real occurrence condition and is independent
of $s$. Both inner products use the original realizations. No local
choice of a meromorphic section at a level place has been inserted.
% END Y87R2

% BEGIN Y87R3
\paragraph{Y87R3. The Euler product for one scalar component.}
Denote by $Z_u(s)$ the doubling integral of this flat section against
the pair of scalar components, with the determinant twist of the
doubling construction. In the convention where the one-variable
matrix coefficient is $\langle R(g)u,u\rangle$, the inducing
character is $\chi_W^{-1}$. The open-orbit identity gives
\[
 Z_u(s)=\int_{U(A)(\mathbb A_F)}
    \langle\omega_B(g)\Psi,\Psi\rangle H(g)^{s-s_*}
                   \langle R(g)u,u\rangle\,dg
\]
on a sufficiently far right half-plane. One can see the orbit
directly because $A$ is definite. A rational maximal isotropic
space in $A\oplus-A$ has zero intersection with either summand,
since $A$ contains no rational isotropic line. Its two projections
are thus isomorphisms. It is the graph of an isometry of $A$, so
the rational isotropic spaces form the single orbit with diagonal
stabilizer. Summing over this orbit unfolds one copy of $[U(A)]$,
where brackets denote the quotient by rational points. Integrating
the other copy gives the matrix coefficient. The defining Eisenstein
series is absolutely convergent sufficiently far right, uniformly
on the compact source quotients. This justifies the sum and
integral interchange. This is the unitary basic doubling identity
of Piatetski-Shapiro and Rallis in the normalization of Harris,
Li and Skinner, \S3.1.2, equation (3.1.2.8).

Here $u$ may have several irreducible summands with the same
unramified Hecke eigenvalues. To remove a place $v\notin S$,
hold all the other components fixed, and normalize the hyperspecial
subgroup $K_v$ to have measure one. The local dependence of the
automorphic coefficient is a bi-$K_v$-invariant function $c_v$.
Right convolution acts through the spherical Hecke character of
$u$. Integrate over $K_v g K_v$ and evaluate at the identity.
This expresses $c_v(g)$ as the corresponding Hecke eigenvalue
times $c_v(1)/\operatorname{vol}(K_vgK_v)$, with inverse cosets
if the inverse convolution convention is used throughout. It also
shows that $c_v$ vanishes identically when $c_v(1)=0$.
The local dependence of the flat section is a scalar multiple of
the spherical section normalized to be one at the identity: the
two local Weil factors have the lattice vacuum on their self-dual
coordinates, and the spherical line of the unramified induction is
one-dimensional by Iwasawa decomposition. More explicitly, at an
odd place invariance under twice the lattice is full lattice
invariance because $2$ is a unit. At a split dyadic place the
two separate split-idempotent clauses give invariance under
translations by the lattice in the first half and modulations by its
annihilator lattice in the conjugate half. The first makes a
Schr\"odinger test constant on lattice cosets, the second forces
its support into the lattice by character orthogonality. The common
line is therefore the characteristic-function line even above $2$.
Doubling tensors the test with its conjugate on the negative
summand. Restriction and negative restriction carry each half into
the same half and preserve both lattice translation conditions.
Thus the doubled vacuum has the self-dual direct-sum lattice, and
its value at the identity supplies the normalization. Thus the spherical formula applies to the
actual partial integrand, even if the global vectors do not factor.

Write $P_w$ for the Satake multiset from $\mathrm{Form}$,
$q_w=\mathrm{N}w$, $z_w=\chi_{W,w}(\varpi_w)$, and let
$\epsilon$ be the quadratic character of $E/F$. The multiplier from
the local integration is
\[
 \frac{\displaystyle\prod_{w\mid v}\prod_{\alpha\in P_w}
              (1-\alpha z_w^{-1}q_w^{-(s+1/2)})^{-1}}
      {\displaystyle\prod_{j=0}^{n-1}
          L_v(2s+n-j,\epsilon_v^{\,b+j})}.
\]
Here we use the split and inert unramified doubling calculation
with hyperspecial volume one, as stated with the corrected
normalizer in Eischen--Harris--Li--Skinner, \S4.2.1, (53)
(\href{https://doi.org/10.1017/fmp.2020.4}{Forum Math. Pi 8 (2020)}).
The restriction of $\chi_W^{-1}$ to $F_v^\times$ is
$\epsilon_v^r=\epsilon_v^b$. At a split place the numerator takes
both components. At an inert place there is one $w$ with $q_w=q_v^2$,
and the $j$-th denominator factor uses
$\epsilon_v(\varpi_v)^{b+j}=(-1)^{b+j}$. The CM host is unramified at
finite primes: its conjugation acts nontrivially on every
conjugation-stable prime's residue field. Dyadic places are split
by $\mathrm{Places}$. These are precisely the unramified cases of
the spherical formula.

Exhaust the adeles by allowing arbitrary local group elements at
finitely many more places outside $S$. On each such region the
preceding calculation removes these additional integrations one at
a time. The unremoved integral is $C$ from Y87R2, for the compatible
Haar measures. Absolute convergence gives the limit of the integrals.
The Euler products converge there as well by the root bounds in
$\mathrm{Form}$ and the finite length of the string. Thus
\[
                    Z_u(s)=C\,L^S(s+1/2)/d^S(s).
\]
The doubled test is $K$-finite: at finite places its stabilizer is open
in the compact group, and at the real places the Gaussian times a
polynomial of bounded degree belongs to a finite sum of compact
types. The standard Eisenstein continuation for this section
makes $Z_u$ meromorphic. Indeed the regularized local holomorphic
family obtained by clearing a pole is locally uniformly bounded
on the compact pair of source quotients, so that integration there
is a continuous functional on this family. By uniqueness the
displayed identity continues through $\Re s>r/2$ to a neighborhood
of $s_*$. Y86N1
and the holomorphic, nonzero normalizer of Y86N2 now give
$\operatorname{Res}_{s=s_*}Z_u(s)\ne0$.
% END Y87R3

% BEGIN Y87R4
\paragraph{Y87R4. Vanishing against every doubled theta kernel.}
For the same $u$ suppose that every downward lift $I_u(f,h)$ to
$U(W)$, as $f$ and $h$ vary, vanishes. On a doubled Schwartz
function for $(A\oplus -A,W)$ define $D_u(\,\cdot\,;h)$
by integration of its theta kernel on $[U(A)]^2$ against
\[
       u(g_1)\overline{u(g_2)}\,\chi_W(\det g_2).
\]
For $\delta(f_1,f_2)$, restriction of the doubled Weil action to
the two factors gives the first action and the conjugate of the
second, multiplied by $\chi_W(\det g_2)^{-1}$. This follows by
writing the second summand with negative form and conjugating the
Heisenberg action there. Its central character reverses; the ratio
of the two splittings on that summand is the displayed determinant
character. Rational lattice points in the direct sum separate into
pairs. The determinant factors cancel before integration and Fubini
gives
\[
       D_u(\delta(f_1,f_2);h)
          = I_u(f_1,h)\,\overline{I_u(f_2,h)}=0.
\]

The vanishing extends to every doubled Schwartz function, for each
fixed $h$. In compatible polarizations the doubling transform is a
linear change of variables and a partial Fourier transform on
Schwartz functions. Products in the two factors are dense. At
the finite places every locally constant compactly supported
function is a finite linear combination of rectangle indicators
on a sufficiently fine product of lattices. At infinity the
Hermite expansion converges in Schwartz seminorms and the Hermite
basis factors in the two sets of variables. The theta functional
is continuous for this approximation: after fixing a finite
lattice support, the rational sum at infinity is bounded by a
Schwartz seminorm times a convergent lattice sum of order greater
than the real dimension. The bound is uniform on compact sets of
translates by continuity of the Weil action on Schwartz space.
The source quotient is compact, so $u$ is bounded and integration
over its square preserves this limit. Equivalently one can transport
the products to the oriented realization by the same Fourier
transform and take compact families of translates there. Thus all
the pairings required in $\mathrm{Kills}$ vanish.
% END Y87R4

% BEGIN Y87R5
\paragraph{Y87R5. The complementary space and the nonvanishing of a varying lift.}
Consider the meromorphic Siegel Eisenstein series on the split
doubled group $U(A\oplus -A)$ attached to the flat section from
Y87R2. The complementary-space first term identity applies to
this particular section. In the notation of Gan--Qiu--Takeda,
Theorem 7.3(i), $d(n)=n$ for a unitary pair. The two target
dimensions are $m=b=r+2t$ and $m'=r>0$, with
$m+m'=2n$ and $m>m'$. They lie in the same Witt tower because
$B=W\perp H$. If $W$ has Witt index $j$, the index of $B$ is
$j+t\le n$. The Siegel--Weil point for $B$ is $(m-n)/2=s_*$.
The inducing character in this convention restricts to
$\epsilon^b=\epsilon^r$ on $\mathbb A_F^\times$, as required for
both target spaces; use the same additive character for the two
Weil actions.

Using W87N2, pass to the diagonal
polarization of $A\oplus-A$. The quotient-lattice realization of
Y87R2 and the Schr\"odinger realization are interchanged by Fourier
expansion along the rational Lagrangian. Its zero coefficient
after applying a group element is precisely the section evaluation
in Y87R2. The inverse transform is periodization, by Poisson
summation with the covolume-one measure. On the Siegel Levi this
transports the action by change of variables with its determinant
character and modulus, on the unipotent radical by multiplication
by the additive character of the Gram form, and on the long Weyl
element by Fourier transform. These operators generate the
doubled group and are the Schr\"odinger Weil action used in the
identity. The same change of polarization for $W$ carries
evaluation of its theta distribution back to the original theta
sum. In a standard unitary chart with skew scalar $\mu$, the character
on the chart is $x\mapsto\psi(\lambda x)$, with
$\lambda=-\xi_0/\mu\in F^\times$, where $\xi_0$ is the host's
fixed skew scalar. The quotient belongs to $F^\times$ because complex
conjugation negates both numerator and denominator. This change in
the additive character compensates the change in the Gram coordinate
and is common to $B$ and $W$. Finally, changing the maximal compact
used to extend the same Siegel section at $s_*$ does not change the
residue. By Proposition 6.1 of Gan--Qiu--Takeda the Eisenstein
operator has at most a simple pole there, and the difference of the
two holomorphic sections is divisible by $s-s_*$.

The identity expresses the Eisenstein residue as a nonzero
constant times the appropriate coefficient of the regularized
theta integral on $U(W)$ applied to
$\operatorname{Ik}(\operatorname{pr}_{K_B}\delta(\Psi,\Psi))$.
Here one first averages over the indicated maximal compact of
$U(B)$; the Ikeda map integrates over the added isotropic variables
in a mixed model. In particular the result is an actual doubled
Schwartz test on $W$. For $j=0$ the coefficient is the ordinary
theta integral, as the quotient of $U(W)$ is compact. For $j>0$
it is coefficient $-1$ of the regularized theta integral.
These are the two unitary cases in
\href{https://arxiv.org/pdf/1207.4709}{Gan--Qiu--Takeda,
\S\S3.5, 6.2 and Theorem 7.3(i)}.

Suppose all downward lifts for $u$ vanish. Apply the pullback
pairing with $u$ on the two copies of $U(A)$ to that first term
identity. In the anisotropic case interchange the two compact
integrations and use Y87R4 at each $h\in U(W)(\mathbb A_F)$.
For $j>0$ use the defining regularization on a right half-plane
away from its scalar poles. Its numerator integrates the theta
kernel of a Schwartz test acted on by the regularizing operator
against an auxiliary Eisenstein series on $U(W)$. The modified
theta kernel is rapidly decreasing in $h$, locally uniformly in
the doubled-group variable; the two source integrations range
over compact sets. This gives absolute convergence and Fubini
against the moderate growth of the auxiliary Eisenstein series.
Y87R4 kills the modified test as well. The pullback of the
regularized family is therefore zero on that half-plane and is
zero meromorphically, including its coefficient of order $-1$.
Extracting the Eisenstein residue commutes with the compact
pullback integral as in Y87R3. In both cases
$\operatorname{Res}_{s=s_*}Z_u(s)=0$, contrary to Y87R3.

This proves the source implication $\mathrm{DownNe}$ under its
stated arithmetic, form and occurrence hypotheses. For each
such scalar component there are a Schwartz function and a target
argument at which its downward integral is nonzero. The test and
argument can depend on the component.
% END Y87R5
